% Vietnamese translation for OI Public Library
% Source-PDF: 模板 TEMPLATES/NOIP2017模板复习手册 - Kvar_ispw17.pdf
\documentclass[11pt]{article}
\usepackage{oipl-vn}
\setmonofont{Noto Sans Mono}[Scale=MatchLowercase]

\title{Sổ tay ôn tập mẫu thuật toán NOIP2017}
\author{Kvar\_ispw17}
\date{9 tháng 11, 2017}

\begin{document}
\maketitle

\tableofcontents

\section{Toán học}

\subsection{Ước chung lớn nhất O(1)}

\begingroup
\footnotesize
\begin{verbatim}
#include <cstdio>
#include <algorithm>
#include <cstring>
#include <cmath>
using namespace std;

typedef long long ll;
const int N = 10000000 + 5;
const int sqrtN = sqrt(N);
int st[N], tot, d[N][5];
int _gcd[sqrtN + 5][sqrtN + 5];
bool notp[N];

void sieve() {
    notp[1] = d[1][0] = d[1][1] = d[1][2] = 1;
    for (ll i = 2; i <= N; i++) {
        if (!notp[i]) {
            d[i][0] = d[i][1] = 1;
            d[i][2] = i;
            st[++tot] = i;
        }
        for (ll j = 1; j <= tot && i * st[j] <= N;
           j++) {
            ll k = i * st[j];
            notp[k] = true;
            ll p = d[i][0] * st[j];
            if (p < d[i][1]) {
                d[k][0] = p;
                d[k][1] = d[i][1];
                d[k][2] = d[i][2];
            } else if (p < d[i][2]) {
                d[k][0] = d[i][1];
                d[k][1] = p;
                d[k][2] = d[i][2];
            } else {
                d[k][0] = d[i][1];
                d[k][1] = d[i][2];


                d[k][2] = p;
            }
            if (i % st[j]); else break;
        }
    }
}

void calc() {
    for (ll i = 0; i <= sqrtN; i++)
        _gcd[i][0] = _gcd[0][i] = i;
    for (ll i = 1; i <= sqrtN; i++) {
        for (ll j = 1; j <= i; j++)
            _gcd[i][j] = _gcd[j][i] = _gcd[i - j][j];
    }
}

ll gcd(ll a, ll b) {
    int *x = d[a], g = 1;
    for (ll i = 0; i < 3; i++) {
        ll d;
        if (x[i] <= sqrtN)
            d = _gcd[x[i]][b % x[i]];
        else if (b % x[i])
            d = 1;
        else d = x[i];
        g *= d; b /= d;
    }
    return g;
}

ll A, B;

int main() {
    sieve();
    calc();
    return 0;
}
\end{verbatim}
\endgroup

\subsection{Thuật toán Euclid mở rộng}

\begingroup
\footnotesize
\begin{verbatim}
int exgcd(int a, int b, int &x, int &y) {
    if (b == 0) { x = 1; y = 0; return a; }
    int ret = exgcd(b, a % b, y, x);
    y -= a / b * x;
    return ret;
}
/**
* other solutions:
* {x} = x + b / gcd(x, y) * t
* {y} = y - a / gcd(x, y) * t
* where t in Z+
*/
\end{verbatim}
\endgroup

\subsection{Nhân nhanh số nguyên dài}

\begingroup
\footnotesize
\begin{verbatim}
#define EPS 1e-8
long long mul(long long a, long long b, long long p) {
    long long ret = a * b - (long long)((long
       double)a / p * b + EPS) * p;
    return ret < 0 ? ret + p : ret;
}

long long mul(long long a, long long b, long long p) {
    long long ret = 0;
    while (b) {
        if (b & 1) ret = (ret + a) % p;
        a = (a + a) % p;
        a >>= 1;
    }
    return ret;
}
\end{verbatim}
\endgroup

\subsection{Sàng Euler}

\begingroup
\footnotesize
\begin{verbatim}
#pragma GCC optimize("O3")
#include <algorithm>
#include <cstdio>
#include <cstring>
#include <cmath>
#define N int(1e7 + 10)
using namespace std;

int num_mindiv[N];
int min_div[N];
int sum_div[N], t1[N], t2[N];
int num_div[N];
int phi[N];
int miu[N];
bool is_prime[N];
int prime[N], tot;

void sieve(int n) {
    miu[1] = 1;
    memset(is_prime, true, sizeof is_prime);
    for (int i = 2; i <= n; i++) {
        if (is_prime[i]) {
            min_div[i] = prime[++tot] = i;
            phi[i] = i - 1;
            num_mindiv[i] = 1;
            num_div[i] = 2;
            miu[i] = -1;
            sum_div[i] = 1 + i;
            t1[i] = 1 + i;
            t2[i] = i;
        }
        for (int j = 1; j <= tot && prime[j] * i <=
           n; j++) {
            int k = prime[j] * i;
            is_prime[k] = false;
            if (i % prime[j] == 0) {
            phi[i * prime[j]] = phi[i] * prime[j];
            num_mindiv[k] = num_mindiv[i] + 1;



           num_div[k] = num_div[i] / (num_mindiv[i]
              + 1) * (num_mindiv[k] + 1);
           miu[i] = 0;
           t2[k] = t2[i] * prime[j];
           t1[k] = t1[i] + t2[k];
           sum_div[k] = sum_div[i] / t1[i] * t1[k];
           break;
        }
        phi[k] = phi[i] * phi[prime[j]];
        num_mindiv[k] = 1;
        num_div[k] = num_div[i] * 2;
        miu[k] = -miu[i];
        sum_div[k] = sum_div[i] * sum_div[prime[j]];
        t1[k] = 1 + prime[j];
        t2[k] = prime[j];
    }
}
}

int main() {
    sieve(10000000);
    return 0;
}
\end{verbatim}
\endgroup

\subsection{Thuật toán BSGS}

\begingroup
\footnotesize
\begin{verbatim}
# include <stdio.h>
# include <math.h>
# include <map>
using namespace std;
typedef long long LL;
LL bsgs(LL A, LL B, LL C) {
    LL m, v, e = 1, i;
    m = ceil(sqrt(C));
    v = pow(A, C - 1 - m, C);
    map<LL, LL> hash;
    hash[1] = m;
    for (i = 1; i < m; ++i) {
        e = (e * A) % C;
        if (!hash[e]) hash[e] = i;
    }
    for (i = 0; i < m; ++i) {
        if (hash[B]) {
            LL ret = hash[B];
            hash.clear();
            return i * m + (ret == m ? 0 : ret);
        }
        B = (B * v) % C;
    }
    return -1;
}
int main() {
    LL A, B, C;
    // A ^ x = B mod C
    while (scanf("%lld%lld%lld", &C, &A, &B) != EOF) {
        LL ans = bsgs(A, B, C);
        if (ans == -1) printf("no solution\n");
        else printf("%lld\n", ans);
    }
    return 0;
}
\end{verbatim}
\endgroup

\subsection{Định lý phần dư Trung Hoa}

\begingroup
\footnotesize
\begin{verbatim}
#define LL long long
LL china(int n, LL *m, LL *a) {
    LL M = 1, d, y, x = 0;
    for (int i = 1; i <= n; i++) M *= m[i];
    for (int i = 1; i <= n; i++) {
        LL w = M / m[i];
        exgcd(m[i], w, d, d, y);
        x = (x + y * w * a[i]) % M;
    }
    return (x + M) % M;
}
LL m[15], a[15];
int main() {
    int n;
    scanf("%d", &n);
    for (int i = 1; i <= n; i++)
    scanf("%lld%lld", &m[i], &a[i]);
    printf("%lld", china(n, m, a));
}
\end{verbatim}
\endgroup

\subsection{Tính nghịch đảo modulo trong O(n)}

\begingroup
\footnotesize
\begin{verbatim}
void get_inv() {
    inv[1] = 1;
    for (int i = 2; i <= n; i++)
    inv[i] = (n - n / i) * inv[n % i] % n;
}
\end{verbatim}
\endgroup

\subsection{Số nguyên độ chính xác cao}

\begingroup
\footnotesize
\begin{verbatim}
/**
*   This is a implement for Big Integer in C++.
*   @author Kvar_ispw17
*   @email enkerewpo@gmail.com
*/

#include <bits/stdc++.h>
using namespace std;

typedef long long LL;

const int N = 20005;
const LL zip = 1e8;

LL Pow(LL a, LL b) { LL ret = 1; while (b) { if (b &
   1) {ret *= a;} a *= a; b >>= 1; } return ret; }

class BigInt {
    public:
    LL a[N];
    bool minus;
    BigInt() {
        minus = false;
        memset(a, 0, sizeof a);
    }
    template <typename _Tp>
    BigInt(_Tp x) {
        int p = 1, tmp = 0, y;
        a[0] = 1;
        while (x) {
            y = x % 10;
            tmp += y * Pow(10, p - 1);
            p %= 8;
            x /= 10;
            if (p == 0) {
                a[a[0]] = tmp;
                tmp = 0;
                a[0]++;
            }


       p++;
   }
   if (tmp) a[a[0]] = tmp;
   else a[0]--;
}
void read() {
    char s[N];
    scanf("%s", s + 1);
    int len = strlen(s + 1), st = 1;
    if (s[1] == '-') st = 2, minus = true;
    int p = 1, tmp = 0, y;
    a[0] = 1;
    for (int i = len; i >= st; i--) {
        y = (s[i] - '0') % 10;
        tmp += y * Pow(10, p - 1), p %= 8;
        if (p == 0) {
            a[a[0]] = tmp;
            tmp = 0;
            a[0]++;
        }
        p++;
    }
    if (tmp) a[a[0]] = tmp;
    else a[0]--;
}
void print() {
    if ((a[0] == 0) ||
       ((a[0] == 1) &&
        (a[1] == 0))) {
        printf("0");
        return;
    }
    if (minus) {
        printf("-");
        if (a[a[0]]) printf("%lld", a[a[0]]);
        for (int i = a[0] - 1; i; i--)
            printf("%08lld", a[i]);
        return;
    }
    if (a[a[0]]) {
        printf("%lld", a[a[0]]);


    }
    for (int i = a[0] - 1; i; i--)
        printf("%08lld", a[i]);
}
BigInt abs() {
    BigInt ret = *this;
    ret.minus = 0;
    return ret;
}
LL operator [] (const unsigned int idx) const {
    return a[idx];
}
LL& operator [] (const unsigned int idx) {
    return a[idx];
}
friend bool operator ==
(const BigInt lhs, const BigInt rhs) {
    if (lhs.minus ^ rhs.minus) return false;
    if (lhs[0] == rhs[0]) {
        for (int i = 1; i <= lhs[0]; i++)
        if (lhs[i] != rhs[i])
            return false;
        return true;
    } else return false;
}
friend bool operator >
(const BigInt lhs, const BigInt rhs) {
    if (lhs[0] == rhs[0]) {
        for (int i = 1; i <= lhs[0]; i++) {
            if (lhs[i] > rhs[i])
                return true;
            else return false;
        }
        return false;
    } else return lhs[0] > rhs[0];
}
friend BigInt operator + (BigInt a, BigInt b) {
    BigInt ret;
    bool f = a.minus ^ b.minus, g = false;
    if (f) {
        if (a.minus) g = true;


       a = a.abs(), b = b.abs();
       ret = a - b;
       if (g) ret.minus ^= 1;
       return ret;
   }
   for (int i = 1; i <= max(a[0], b[0]); i++) {
       ret[i] += a[i] + b[i];
       ret[i + 1] = ret[i] / zip;
       ret[i] %= zip;
   }
   ret[0] = max(a[0], b[0]);
   if (ret[ret[0] + 1]) {
       ret[0]++;
   }
   if (a.minus) ret.minus = true;
   return ret;
}
BigInt operator * (int rhs) {
    bool f = (rhs < 0) ^ minus;
    BigInt c;
    rhs = std::abs(rhs);
    for (int i = 1; i <= a[0]; i++) {
        c.a[i] += a[i] * rhs;
        if (c.a[i] >= zip) {
            c.a[i + 1] += c.a[i] / zip;
            c.a[i] %= zip;
        }
    }
    c.a[0] = a[0];
    if (c.a[c.a[0] + 1] > 0) c.a[0]++;
    if (f) c.minus = true;
    return c;
}
BigInt operator * (const BigInt rhs) const {
    BigInt ret;
    bool f = false;
    if (minus ^ rhs.minus) f = true;
    for (int i = 1; i <= a[0]; i++)
    for (int j = 1; j <= rhs[0]; j++) {
        ret[i + j - 1] += a[i] * rhs[j];
        ret[i + j] += ret[i + j - 1] / zip;


       ret[i + j - 1] %= zip;
   }
   ret[0] = max(0LL, a[0] + rhs[0] - 1);
   if (ret[a[0] + rhs[0]]) {
       ret[0]++;
   }
   if (f) ret.minus = true;
   return ret;
}
friend BigInt operator -
(BigInt lhs, BigInt rhs) {
    if (lhs == rhs) {
        return BigInt(0);
    }
    if (lhs.minus ^ rhs.minus) {
        if (rhs.minus) return lhs + rhs.abs();
        else {
            BigInt ret = lhs.abs() + rhs.abs();
            ret.minus = true;
            return ret;
        }
    }
    bool right = lhs > rhs;
    lhs = lhs.abs(), rhs = rhs.abs();
    if (rhs > lhs) swap(lhs, rhs);
    BigInt ret;
    memcpy(ret.a, lhs.a, sizeof lhs.a);
    for (int i = 1; i <= lhs.a[0]; i++) {
        ret[i] -= rhs[i];
        if (ret[i] < 0)
            ret[i + 1]--, ret[i] += zip;
    }
    while (ret[0] >= 1 && !ret[ret[0]])
        ret[0]--;
    if (!right) ret.minus = true;
    return ret;
}
BigInt operator / (int rhs) {
    BigInt c;
    LL x = 0;
    for (int i = a[0]; i; i--) {


             x = x * zip + a[i];
             c.a[i] = x / rhs, x = x % rhs;
        }
        c.a[0] = a[0];
        if (c.a[0] && !c.a[c.a[0]])
            c.a[0]--;
        return c;
     }
     LL operator % (int rhs) {
     BigInt c;
     LL x = 0;
     for (int i = a[0]; i; i--) {
         x = x * zip + a[i];
         c.a[i] = x / rhs;
         x = x % rhs;
     }
     c.a[0] = a[0];
     if (c.a[0] && !c.a[c.a[0]])
         c.a[0]--;
     return x;
}
};

int main() {
    BigInt a, b;
    a.read(), b.read();
    BigInt c = a + b;
    c.print(); puts("");
    c = a - b;
    c.print(); puts("");
    c = a * b;
    c.print(); puts("");
    return 0;
}
\end{verbatim}
\endgroup

\subsection{Biến đổi Fourier nhanh}

\begingroup
\footnotesize
\begin{verbatim}
#include <bitset>
#include <cstdio>
#include <cstring>
#include <cmath>
#include <algorithm>
#define N 400005
#define pi acos (-1.0) // PI value
using namespace std;
struct complex {
    double r, i;
    complex (double real = 0.0, double image = 0.0) {
        r = real; i = image;
    }
    // The following is the definition of three kinds
       of imaginary arithmetic
    complex operator + (const complex o) {
        return complex (r + o.r, i + o.i);
    }
    complex operator - (const complex o) {
        return complex (r - o.r, i - o.i);
    }
    complex operator * (const complex o) {
        return complex (r * o.r - i * o.i, r * o.i +
           i * o.r);
    }
} x1 [N], x2 [N];
char a [N / 2], b [N / 2];
int sum [N]; // The result exists in sum
int vis [N];
void brc (complex * a, int l) {
    // original bingliang binary parallel reversal
       replacement Taifeng do not understand, wrote
       an O (n) ......
    memset (vis, 0, sizeof (vis)); // O (logn) is
       followed
    for (int i = 1; i < l - 1; i ++) {
        int x = i, y = 0;
        int m = (int) log2 (l) + 0.1;
        if (vis [x]) continue;


        while (m--) {
            y <<= 1;
            y |= (x & 1);
            x >>= 1;
        }
        vis [i] = vis [y] = 1;
        swap (a [i], a [y]);
    }
}
void fft (complex * y, int l, double on) {// FFT O
   (nlogn)
    // where DFT is on== 1, on== - 1 is IDFT
    register int h, i, j, k;
    complex u, t;
    brc (y, l); // Call reversal permutation
    for (h = 2; h <= l; h <<= 1) { // Control the
       number of layers
        // initial unit complex root
        complex wn (cos (on * 2 * pi / h), sin (on *
           2 * pi / h));
        for (j = 0; j < l; j += h) { // Controls the
           start subscript
            complex w (1, 0); // Initialize the
               spiral factor
            for (k = j; k < j + h / 2; k ++) { //
               paired
                u = y [k];
                t = w * y [k + h / 2];
                y [k] = u + t;
                y [k + h / 2] = u - t;
                w = w * wn; // Update the spiral
                   factor
            } // It is said that the operation above
               is called butterfly operation ...
        }
    }
    if (on == -1) for (i = 0; i < l; i ++) y [i] .r
       /= l; // IDFT
}
int main () {
    int l1, l2, l;


    register int i;
    while (scanf("%s%s", a, b) != EOF) {
        l1 = strlen (a);
        l2 = strlen (b);
        l = 1;
        while (l < l1 * 2 || l < l2 * 2) l <<= 1;
           //The number of boundaries into 2 ^ n
        // with the dichotomy and reverse replacement
        for (i = 0; i < l1; i ++) { // Inverted
            x1 [i] .r = a [l1 - i - 1] - '0';
            x1 [i] .i = 0.0;
        }
        for (; i < l; i ++) x1 [i] .r = x1 [i] .i =
           0.0;
        // Initialize the extra count to zero
        for (i = 0; i < l2; i ++) {
            x2 [i] .r = b [l2 - i - 1] - '0';
            x2 [i] .i = 0.0;
        }
        for (; i < l; i ++) x2 [i] .r = x2 [i] .i =
           0.0;
        fft (x1, l, 1); // DFT (a)
        fft (x2, l, 1); // DFT (b)
        for (i = 0; i < l; i ++) x1 [i] = x1 [i] * x2
           [i];
        fft (x1, l, -1); // IDFT (a * b)
        for (i = 0; i < l; i ++) sum [i] = x1 [i] .r
           + 0.5; // round off
        for (i = 0; i < l; i ++) { // Carry
            sum [i + 1] += sum [i] / 10;
            sum [i] %= 10;
        }
        l = l1 + l2 - 1;
        while (sum [l] <= 0 && l > 0) l--; //
           retrieve the most significant bit
        for (i = l; i >= 0; i--) putchar (sum [i] +
           '0'); // output in reverse
        putchar ('\n');
    }
    return 0;
}
\end{verbatim}
\endgroup

\subsection{Kiểm tra nguyên tố Miller--Rabin}

\begingroup
\footnotesize
\begin{verbatim}
#include <bits/stdc++.h>

/* QUICK POWER */
template <typename T>
T power(T A, T B, T P) {
    T ret = 1;
    while (B) {
        if (B & 1) ret = ret * A % P;
        A = A * A % P; B >>= 1;
    } return ret;
}

/* MILLER RABIN PRIME */
const int u[15] = {2, 3, 5, 7, 11, 13, 17, 19, 23,
   29, 31, 37}, cas = 12;

template <typename T>
bool miller_check(T a, T p) {
    T x = p - 1;
    if (power<T>(a, p - 1, p) != 1) return false;
    while (~x & 1) x >>= 1; /* x is an even number */
    a = power<T>(a, x, p);
    while (x < p - 1) {
        if (a != 1 && a != p - 1) {
            a = a * a % p;
            if (a == 1) return false;
        } else a = a * a % p;
        x <<= 1;
    }
    return true;
}
template <typename T>
bool miller(T p) {
    if (p == 1 || p == 0) return 0;
    for (int i = 0; i < cas; i++) {
        //printf("P: %d NOW: %d\n", p, u[i]);
        if (p == u[i]) return true;
        if (!miller_check<T>(u[i], p)) return false;
    }


    return true;
}

int main() {
    return 0;
}
\end{verbatim}
\endgroup

\subsection{Phân tích thừa số Pollard Rho}

\begingroup
\footnotesize
\begin{verbatim}
#include <bits/stdc++.h>
/**
* This is a Template of Pollard Rho and .Miller Rabin
   Prime Test
* email: enkerewpo@gmail.com
* blog: enkerewpo.github.io
*/
ll mult_mod(ll a, ll b, ll c) {
    a %= c;
    b %= c;
    ll ret = 0;
    while (b) {
        if (b & 1) {
            ret += a;
            ret %= c;
        }
        a <<= 1;
        if (a >= c)a %= c;
        b >>= 1;
    }
    return ret;
}

ll q_pow(ll x, ll n, ll mod) {
    if (n == 1) return x % mod;
    x %= mod;
    ll tmp = x;
    ll ret = 1;
    while (n) {
        if (n & 1) ret = mult_mod(ret, tmp, mod);
        tmp = mult_mod(tmp, tmp, mod);


        n >>= 1;
    }
    return ret;
}

bool check(ll a, ll n, ll x, ll t) {
    ll ret = q_pow(a, x, n);
    ll last = ret;
    for (int i = 1; i <= t; i++) {
        ret = mult_mod(ret, ret, n);
        if (ret == 1 && last != 1 && last != n - 1)
           return true;
        last = ret;
    }
    if (ret != 1) return true;
    return false;
}

bool Miller_Rabin(ll n) {
    if (n < 2) return false;
    if (n == 2) return true;
    if ((n & 1) == 0) return false;
    ll x = n - 1;
    ll t = 0;
    while ((x & 1) == 0) {
        x >>= 1;
        t++;
    }
    for (int i = 0; i < S; i++) {
        ll a = rand() % (n - 1) + 1;
        if (check(a, n, x, t))
        return false;
    }
    return true;
}

ll fac[100];
int tot;

ll gcd(ll a, ll b) {
    if (a == 0) return 1;


    if (a < 0) return gcd(-a, b);
    while (b) {
        ll t = a % b;
        a = b;
        b = t;
    }
    return a;
}

ll Pollard_rho(ll x, ll c) {
    ll i = 1, k = 2;
    ll x0 = rand() % x;
    ll y = x0;
    while (1) {
        i++;
        x0 = (mult_mod(x0, x0, x) + c) % x;
        ll d = gcd(y - x0, x);
        if (d != 1 && d != x) return d;
        if (y == x0) return x;
        if (i == k) {
            y = x0;
            k += k;
        }
    }
}

void findfac(ll n) {
    if (Miller_Rabin(n)) {
        fac[tot++] = n;
        return;
    }
    ll p = n;
    while (p >= n)
    p = Pollard_rho(p, rand() % (n - 1) + 1);
    findfac(p);
    findfac(n / p);
}
\end{verbatim}
\endgroup

\section{Chuỗi}

\subsection{Tự động cơ AC}

\begingroup
\footnotesize
\begin{verbatim}
#include <cstdio>
#include <string>
#include <cstring>
#include <algorithm>
using namespace std;
const int maxn = 100000;
struct point {
    int nxt[26];
    int fail, cnt;
} a[maxn];

int cnt, root;
int que[maxn], head, tail;
char A[maxn];

inline void insert(char* s) {
    int now = root;
    for (int i = 0; s[i] != '\0'; ++i) {
        int x = s[i] - 'a';
        if (a[now].nxt[x] == 0) a[now].nxt[x] = cnt++;
        now = a[now].nxt[x];
    }
    return;
}

inline void build() {
    for (int i = 0; i < 26; ++i) {
        if (a[root].nxt[i] != 0) {
            int p = a[root].nxt[i];
            a[p].fail = root;
            que[tail++] = p;
        }
    }
    while (head < tail) {
        int p = que[head++];




       for (int i = 0; i < 26; ++i) if (a[p].nxt[i])
          {
            int x = a[p].nxt[i];
            int j = a[p].fail;
            while (j != root && a[j].nxt[i] == 0) j =
               a[j].fail;
            if (a[j].nxt[i]) j = a[j].nxt[i];
            a[x].fail = j;
            que[tail++] = x;
       }
    }
    return;
}

inline void solve(char* s) {
    int now = root;
    for (int i = 0; s[i] != '\0'; ++i) {
        int x = s[i] - 'a';
        while (now != root && a[now].nxt[x] == 0) now
           = a[now].fail;
        if (a[now].nxt[x]) now = a[now].nxt[x];
        ++a[now].cnt;
    }
    for (int i = tail - 1; i >= 0; --i) {
        int p = a[que[i]].fail;
        a[p].cnt += a[que[i]].cnt;
    }
    return;
}
int main() {
    int n;
    cnt = 2;
    root = 1;
    scanf("%d", &n);
    while (n--) scanf("%s", A), insert(A);
    build();
    scanf("%s", A);
    solve(A);
    printf("%d", a[2].cnt);
    return 0;
}
\end{verbatim}
\endgroup

\subsection{Thuật toán KMP}

\begingroup
\footnotesize
\begin{verbatim}
#include <bits/stdc++.h>
#define N 10000000 + 5
using namespace std;

typedef long long ll;
ll ans;
char str[N], t[N];

ll nxt[N];
int main() {
    scanf("%s", t + 1);
    scanf("%s", str + 1);
    int lent = strlen(t + 1);
    int lenstr = strlen(str + 1);
    int j = 0;
    nxt[1] = 0;
    for (int i = 2; i <= lent; i++) {
        while (j > 0 && t[i] != t[j + 1]) j = nxt[j];
        if (t[j + 1] == t[i]) j++;
        nxt[i] = j;
    }
    j = 0;
    for (int i = 1; i <= lenstr; i++) {
        while (j > 0 && str[i] != t[j + 1]) j =
           nxt[j];
        if (t[j + 1] == str[i]) j++;
        if (j == lent) {

            ans++;
            j = nxt[j];
        }
    }
    printf("%lld\n", ans);
    return 0;
}
\end{verbatim}
\endgroup

\subsection{Thuật toán Manacher}

\begingroup
\footnotesize
\begin{verbatim}
#include<cstdio>
#include<cstring>
#include<iostream>
#include<algorithm>

using namespace std;
const int maxn = 1000010;

int n, len[maxn * 2];
char S[maxn * 2];

inline void solve() {
    n = strlen(S + 1);
    n = n * 2 - 1;
    for (int i = n; i > 0; --i) if (i & 1) S[i] =
       S[(i + 1) / 2];
    else S[i] = '#';
    int mx = 0, id = 0;
    for (int i = 1; i <= n; ++i) {
        if (mx <= i) len[i] = 1;
        else len[i] = min(mx - i + 1, len[2 * id -
           i]);
        while (i - len[i] > 0 && i + len[i] <= n &&
           S[i + len[i]] == S[i - len[i]]) ++len[i];
        if (i + len[i] - 1 > mx) {
            mx = i + len[i] - 1;
            id = i;
        }
    }
}

int main() {
    scanf("%s", S + 1);
    solve();
    printf("S= %s\n", S + 1);
    for (int i = 1; i <= n; ++i)
       printf("len[%d]=%d\n", i, len[i]);
    return 0;
}
\end{verbatim}
\endgroup

\section{Lý thuyết đồ thị}

\subsection{Phân rã cây theo chuỗi nặng}

\begingroup
\footnotesize
\begin{verbatim}
#include <bits/stdc++.h>
using namespace std;

#define N 30000 + 5
#define inf 0x7f7f7f7f

int n, q, hd[N], nxt[N], to[N], w[N], edge, tot;
int siz[N], dep[N], u_son[N], fa[N], top[N], tr[N],
   dtr[N];

namespace segtree {

   int TR[4 * N], MAX[4 * N];

   void up(int p) {
       TR[p] = TR[p << 1] + TR[p << 1 | 1];
       MAX[p] = max(MAX[p << 1], MAX[p << 1 | 1]);
   }

   void build(int l, int r, int p) {
       if (l == r) {
           TR[p] = w[dtr[l]];
           MAX[p] = TR[p];
           return;
       }
       int mid = (l + r) / 2;
       build(l, mid , p << 1);
       build(mid + 1, r, p << 1 | 1);
       up(p);
   }

   void change(int l, int r, int P, int val, int p) {
       if (l == r && l == P) {
           TR[p] = val;
           MAX[p] = TR[p];
           return;



        }
        int mid = (l + r) / 2;
        if (P <= mid) change(l, mid, P, val, p << 1);
        else change(mid + 1, r, P, val, p << 1 | 1);
        up(p);
    }

    int ask_max(int l, int r, int L, int R, int p) {
        if (l > R || r < L) return -inf;
        if (L <= l && r <= R) return MAX[p];
        int mid = (l + r) / 2;
        return max(ask_max(l, mid, L, R, p << 1),
           ask_max(mid + 1, r, L, R, p << 1 | 1));
    }

    int ask_sum(int l, int r, int L, int R, int p) {
        if (l > R || r < L) return 0;
        if (L <= l && r <= R) return TR[p];
        int mid = (l + r) / 2;
        return ask_sum(l, mid, L, R, p << 1) +
           ask_sum(mid + 1, r, L, R, p << 1 | 1);
    }

}

namespace HLD {

    void add(int a, int b) {
        nxt[++edge] = hd[a], to[edge] = b, hd[a] =
           edge;
        nxt[++edge] = hd[b], to[edge] = a, hd[b] =
           edge;
    }

    void dfs1(int u, int __fa, int __dep) {
        fa[u] = __fa;
        dep[u] = __dep;
        siz[u] = 1;
        for (int c = hd[u]; c; c = nxt[c]) {
            int v = to[c];
            if (v != __fa) {


            dfs1(v, u, __dep + 1);
            siz[u] += siz[v];
            if (!u_son[u] || siz[v] >
               siz[u_son[u]]) u_son[u] = v;
        }
    }
}

void dfs2(int u, int id) {
    top[u] = id;
    tr[u] = ++tot;
    dtr[tot] = u;
    if (!u_son[u]) return;
    dfs2(u_son[u], id);
    for (int c = hd[u]; c; c = nxt[c]) {
        int v = to[c];
        if (v != u_son[u] && v != fa[u]) dfs2(v,
           v);
    }
}

int max(int a, int b) {
    int f1, f2, ret = -inf;
    f1 = top[a];
    f2 = top[b];
    while (f1 != f2) {
        if (dep[f1] < dep[f2]) swap(a, b),
           swap(f1, f2);
        ret = std::max(ret, segtree::ask_max(1,
           n, tr[f1], tr[a], 1));
        a = fa[f1];
        f1 = top[a];
    }
    if (dep[a] < dep[b]) swap(a, b);
    ret = std::max(ret, segtree::ask_max(1, n,
       tr[b], tr[a], 1));
    return ret;
}

int sum(int a, int b) {
    int f1, f2, ret = 0;


        f1 = top[a];
        f2 = top[b];
        while (f1 != f2) {
            if (dep[f1] < dep[f2]) swap(a, b),
               swap(f1, f2);
            ret += segtree::ask_sum(1, n, tr[f1],
               tr[a], 1);
            a = fa[f1];
            f1 = top[a];
        }
        if (dep[a] < dep[b]) swap(a, b);
        ret += segtree::ask_sum(1, n, tr[b], tr[a],
           1);
        return ret;
    }

}

int main() {
    scanf("%d", &n);
    for (int i = 1; i < n; i++) {
        int a, b;
        scanf("%d%d", &a, &b);
        HLD::add(a, b);
    }
    HLD::dfs1(1, 0, 1);
    HLD::dfs2(1, 1);
    for (int i = 1; i <= n; i++) scanf("%d", &w[i]);
    segtree::build(1, n, 1);
    scanf("%d", &q);
    while (q--) {
        string op;
        cin >> op;
        int a, b;
        scanf("%d%d", &a, &b);
        if      (op == "QMAX")     printf("%d\n",
           HLD::max(a, b));
        else if (op == "QSUM")     printf("%d\n",
           HLD::sum(a, b));
        else                       segtree::change(1,
           n, tr[a], b, 1);


    }
    return 0;
}
\end{verbatim}
\endgroup

\subsection{Bao đóng bắc cầu Floyd}

\begingroup
\footnotesize
\begin{verbatim}
void floyd() {
    int i, j, k;
    int n;
    bool* con = new bool[n + 1][n + 1];
    for (int k = 1; k <= n; k++)
    for (int i = 1; i <= n; i++)
    for (int j = 1; j <= n; j++)
    con[i][j] |= con[i][k] & con[k][j];
}
\end{verbatim}
\endgroup

\subsection{Thành phần song liên thông đỉnh}

\begingroup
\footnotesize
\begin{verbatim}
#include <iostream>
#include <cstdio>
#include <algorithm>
#include <cstring>
#include <string>
#include <cmath>
#include <vector>
#include <stack>
using namespace std;

const int maxn = 1000;

struct Edge {// structure of the side of the stack
    int u, v;
    Edge(int uu, int vv) {
        u = uu;
        v = vv;
    }
};
stack <Edge> s;




struct edge {// The edge structure of the chained
   forward starburst
    int v, next;
} edges[maxn];

int n, m; // The number of nodes, the number of
   undirected edges
int e, head[maxn];
int pre[maxn]; // The timestamp of the first visit
int dfs_clock; // timestamp
int iscut[maxn]; // mark whether the node is cut top
int bcc_cnt; // dot_ The number of double connected
   components
int bccno[maxn]; // The point to which the node
   belongs _ The number of the dual connected
   component
vector <int> bcc[maxn]; // dot_double connected
   component

void addedges(int u, int v) {// Add edge
    edges[e].v = v;
    edges[e].next = head[u];
    head[u] = e ++;
    edges[e].v = u;
    edges[e].next = head[v];
    head[v] = e ++;
}

int dfs(int u, int fa) {
    int lowu = pre[u] = ++ dfs_clock;
    int child = 0;
    for (int i = head[u]; i != -1; i = edges[i].next)
       {
         int v = edges[i].v;
         Edge e = (Edge) {u, v};
         if (! pre[v]) {
             s.push(e);
             child ++;
             int lowv = dfs(v, u);
             lowu = min(lowu, lowv); // Update lowu
                with descendants


           if (lowv >= pre[u]) { // Find a subtree
              that satisfies the condition of the cut
               iscut[u] = 1;
               bcc_cnt ++;
               bcc[bcc_cnt].clear();
               for (;;) { // save the bcc information
                   Edge x = s.top(); s.pop();
                   if (bccno[x.u] != bcc_cnt)
                      {bcc[bcc_cnt].push_back(x.u);
                      bccno[x.u] = bcc_cnt;}
                   if (bccno[x.v] != bcc_cnt)
                      {bcc[bcc_cnt].push_back(x.v);
                      bccno[x.v] = bcc_cnt;}
                   if (x.u == u && x.v == v) break;

                }
            }
        } else if (pre[v] < pre[u] && v != fa) { //
           Update lowu
            s.push(e);
            lowu = min(lowu, pre[v]);
        }
    }
    if (fa < 0 && child == 1) iscut[u] = 0; // If the
       root node has only one subtree then it is not
       a cut
    return lowu;
}

void init() {
    memset(pre, 0, sizeof(pre));
    memset(iscut, 0, sizeof(iscut));
    memset(head, -1, sizeof(head));
    memset(bccno, 0, sizeof(bccno));
    e = 0; dfs_clock = 0; bcc_cnt = 0;
}

int main() {
    int u, v;
    while (scanf("%d%d", &n, &m) != EOF) {
        init();


        for (int i = 0; i < m; i ++) {
            scanf("%d%d", &u, &v);
            addedges(u, v);
        }
        dfs(1, -1);
        for (int i = 1; i <= bcc_cnt; i ++) {
            for (int j = 0; j < (int)bcc[i].size(); j
               ++)
                cout << bcc[i][j] << "";
            cout << endl;
        }
    }
    return 0;
}
\end{verbatim}
\endgroup

\subsection{Thành phần song liên thông cạnh}

\begingroup
\footnotesize
\begin{verbatim}
#include <iostream>
#include <cstdio>
#include <algorithm>
#include <cstring>
#include <string>
#include <cmath>
#include <vector>
using namespace std;

const int maxn = 1000;
struct Edge {
    int no, v, next; // no: Number of the edge
} edges[maxn];

int n, m, ebcnum; // number of nodes, number of
   undirected edges, number of edges_double connected
   components
int e, head[maxn];
int pre[maxn]; // The timestamp of the first visit
int dfs_clock; // timestamp
int isbridge[maxn]; // mark if the edge is a bridge
vector <int> ebc[maxn]; // Edge_double connected
   component


void addedges(int num, int u, int v) {// Add edge
    edges[e].no = num;
    edges[e].v = v;
    edges[e].next = head[u];
    head[u] = e++;
    edges[e].no = num++;
    edges[e].v = u;
    edges[e].next = head[v];
    head[v] = e++;
}

int dfs_findbridge(int u, int fa) {// Find all the
   bridges
    int lowu = pre[u] = ++dfs_clock;
    for (int i = head[u]; i != -1; i = edges[i].next)
       {
         int v = edges[i].v;
         if (! pre[v]) {
             int lowv = dfs_findbridge(v, u);
             lowu = min(lowu, lowv);
             if (lowv > pre[u]) {
                 isbridge[edges[i].no] = 1; // bridge
             }
         } else if (pre[v] < pre[u] && v != fa) {
             lowu = min(lowu, pre[v]);
         }
    }
    return lowu;
}

void dfs_coutbridge(int u, int fa) {// Saves the
   information of the edge_double connected component
    ebc[ebcnum].push_back(u);
    pre[u] = ++dfs_clock;
    for (int i = head[u]; i != -1; i = edges[i].next)
       {
         int v = edges[i].v;
         if (! isbridge[edges[i].no] && ! pre[v])
            dfs_coutbridge(v, u);
    }


}

void init() {
    memset(pre, 0, sizeof(pre));
    memset(isbridge, 0, sizeof(isbridge));
    memset(head, -1, sizeof(head));
    e = 0; ebcnum = 0;
}

int main() {
    int u, v;
    while (scanf("%d%d", &n, &m) != EOF) {
        init();
        for (int i = 0; i < m; i++) {
            scanf("%d%d", &u, &v);
            addedges(i, u, v);
        }
        dfs_findbridge(1, -1);
        memset(pre, 0, sizeof(pre));
        for (int i = 1; i <= n; i++) {
            if (! pre[i]) {
                ebc[ebcnum].clear();
                dfs_coutbridge(i, -1);
                ebcnum++;
            }
        }
        for (int i = 0; i < ebcnum; i++) {
            for (int j = 0; j < (int)ebc[i].size();
               j++)
            cout << ebc[i][j] << " ";
            cout << endl;
        }
    }
    return 0;
}
\end{verbatim}
\endgroup

\subsection{Tổ tiên chung gần nhất}

\begingroup
\footnotesize
\begin{verbatim}
#include <bits/stdc++.h>
using namespace std;


int dep[10], fa[10][10];
int lca(int x, int y) {
    if (dep[x] < dep[y]) swap(x, y);
    int d = dep[x] - dep[y];
    for (int i = 0; i <= 18; i++) {
        if (d & (1 << i)) x = fa[x][i];
    }
    for (int i = 18; i >= 0; i--) {
        if (fa[x][i] != fa[y][i]) {
            x = fa[x][i], y = fa[y][i];
        }
    }
    if (x == y) return x;
    else return fa[x][0];
}
\end{verbatim}
\endgroup

\subsection{Hệ ràng buộc hiệu}

\begingroup
\footnotesize
\begin{verbatim}
#include <algorithm>
#include <cstring>
#include <cstdio>
#include <queue>
#include <cassert>

#define N 200005 + 5
#define inf 1e9
using namespace std;
#define int long long

int hd[2 * N], nxt[2 * N], to[2 * N], w[2 * N], tot;
int n, k;
void add(int a, int b, int c) {
    nxt[++tot] = hd[a], to[hd[a] = tot] = b, w[tot] =
       c;
}

int dis[4 * N];
int cnt[4 * N];
bool vis[4 * N];



int spfa() {
    memset(dis, -inf, sizeof dis);
    queue<int> q;
    q.push(0);
    vis[0] = true;
    dis[0] = 0;
    cnt[0]++;
    while (!q.empty()) {
        int u = q.front();
        q.pop();
        if (cnt[u] == n + 1) return false;
        vis[u] = false;
        for (int e = hd[u]; e != -1; e = nxt[e]) {
            int v = to[e];
            if (dis[u] + w[e] > dis[v]) {
                dis[v] = dis[u] + w[e];
                if (!vis[v]) {
                    q.push(v);
                    cnt[v]++;
                    vis[v] = true;
                }
            }
        }
    }
    return true;
}

signed main() {
    memset(hd, -1, sizeof hd);
    scanf("%lld%lld", &n, &k);
    while (k--) {
        int x, a, b;
        scanf("%lld%lld%lld", &x, &a, &b);
        switch (x) {
            case 1 : { add(a, b, 0), add(b, a, 0);
               break;}                           //
               A   =   B
            case 2 : { if (a == b) return puts("-1"),
               0; else add(a, b, 1); break;} // A
               <   B



              case 3 : { add(b, a, 0); break;}
                                                 //     A
                   >= B
              case 4 : { if (a == b) return puts("-1"),
                 0; else add(b, a, 1); break;} // A
                 >   B
              case 5 : { add(a, b, 0); break;}
                                                  // A
                   <= B
          }
      }
      for (int i = n; i >= 1; i--) add(0, i, 1);
      if (spfa()) {
          int ans = 0;
          for (int i = 1; i <= n; i++) ans += dis[i];
          printf("%lld\n", ans);
      } else puts("-1");
      return 0;
}
\end{verbatim}
\endgroup

\subsection{Cây ảo}

\begingroup
\footnotesize
\begin{verbatim}
#include <algorithm>
#include <cstring>
#include <cstdio>
/**
* This is a Hollow Tree implement by Kvar_ispw17
* @author Kvar_ispw17
* @email enkerewpo@gmail.com
*/
using namespace std;

const int N = 10;
const int M = 10;

int n, m, key[5];
int dep[N], fa[N][10], dfn[N], tim;

/**



*   Sample input : |V| = 6, |E| = 5 E = { (1, 2) (2,
    5) (2, 3) (3, 4) (3, 6) }
*   Sample key-node input : V' = { 3, 6, 5 }
*
* Defination : node '0' as the super root.
*           node '1' as the tree root.
*/

int hd[N], nxt[2 * M], to[2 * M], tot;
/**
* add edge to the original tree
* @param a from
* @param b to
*/
void add(int a, int b) {
    nxt[++tot] = hd[a], to[hd[a] = tot] = b;
}

int hd2[N], nxt2[2 * M], to2[2 * M], tot2, from2[2 *
   M];
/**
* add edge to the hollow tree
* @param a from
* @param b to
*/
void add2(int a, int b) {
    nxt2[++tot2] = hd2[a], to2[hd2[a] = tot2] = b,
       from2[tot2] = a;
}

int lca(int a, int b) {
    if (dep[a] < dep[b]) swap(a, b);
    int d = dep[a] - dep[b];
    for (int j = 0; j <= 5; j++) {
        if (d & (1 << j)) a = fa[a][j];
    }
    if (a == b) return a;
    for (int j = 5; j >= 0; j--) {
        if (fa[a][j] != fa[b][j]) {
            a = fa[a][j], b = fa[b][j];
        }


    }
    return fa[a][0];
}

bool Cmp_Fn(const int &a, const int &b) {
    return dfn[a] < dfn[b];
}

/**
* main algorithm for building the hollow tree.
*/
void build() {
    sort(key + 1, key + m + 1, Cmp_Fn);
    int s[N], top = -1;
    memset(s, 0, sizeof s);
    s[++top] = 0;
    int cnt = m;
    for (int i = 1; i <= m; i++) {
        int u = key[i];
        int lca = ::lca(u, s[top]);
        if (lca == s[top]) s[++top] = u;
        else {
            while (top - 1 >= 0 && dep[s[top - 1]] >=
               dep[lca]) {
                add2(s[top - 1], s[top]);
                add2(s[top], s[top - 1]);
                top--;
            }
            if (s[top] != lca) {
                add2(lca, s[top]);
                add2(s[top], lca);
                s[top] = lca;
                key[++cnt] = lca;
            }
            s[++top] = u;
        }
    }
    for (int i = 0; i < top; i++) {
        add2(s[i], s[i + 1]);
        add2(s[i + 1], s[i]);
    }


}

/**
* print the hollow tree int depth by first-order dfs
* @param u current node
* @param d current depth
*/
void print(int u, int _fa, int d) {
    for (int i = 1; i <= d; i++) printf("    ");
    printf("%d\n", u);
    for (int e = hd2[u]; e; e = nxt2[e]) {
        int v = to2[e];
        if (v != _fa) {
            print(v, u, d + 1);
        }
    }
}

void dfs(int u, int _fa, int d) {
    dep[u] = d;
    dfn[u] = ++tim;
    fa[u][0] = _fa;
    for (int e = hd[u]; e; e = nxt[e]) {
        int v = to[e];
        if (v != _fa) {
            dfs(v, u, d + 1);
        }
    }
}

/**
* program main function
* @return program stantard return code
*/
int main() {
    scanf("%d", &n);
    for (int i = 1; i < n; i++) {
        int a, b;
        scanf("%d%d", &a, &b);
        add(a, b);
        add(b, a);


   }
   dfs(1, 0, 1);
   for (int j = 1; j <= 5; j++) {
       for (int i = 1; i <= n; i++) {
           fa[i][j] = fa[fa[i][j - 1]][j - 1];
       }
   }
   scanf("%d", &m);
   for (int i = 1; i <= m; i++) scanf("%d", &key[i]);
   build();
   puts("printing the hollow tree of input by 
      depth:");
   print(0, -1, 0);
   return 0;
}
/*
Sample input:
6
1 2
2 3
2 5
3 4
3 6
3
3 6 4
*/
\end{verbatim}
\endgroup

\subsection{Luồng mạng Dinic}

\begingroup
\footnotesize
\begin{verbatim}
#include <bits/stdc++.h>
using namespace std;
#define N 1000 + 5
#define INF 0x7f7f7f7f
#define M 100000 + 5
int hd[N], nxt[2 * M], f[2 * M], to[2 * M], tot = 1;

int S, T;

void add(int a, int b, int c) {



    nxt[++tot] = hd[a], to[hd[a] = tot] = b, f[tot] =
       c;
    nxt[++tot] = hd[b], to[hd[b] = tot] = a, f[tot] =
       0;
}

int dis[N];
int que[N];

bool check() {
    memset(dis, -1, sizeof dis);
    int head = 0, tail = -1;
    dis[que[++tail] = S] = 0;
    while (head <= tail) {
        int u = que[head++];
        for (int e = hd[u]; e; e = nxt[e]) {
            int v = to[e];
            if (dis[v] == -1 && f[e])
            dis[que[++tail] = v] = dis[u] + 1;
        }
    }
    return dis[T] != -1;
}

int argument(int u, int rem) {
    if (u == T) return rem;
    int cnt = 0;
    for (int e = hd[u]; e && rem > cnt; e = nxt[e]) {
        int v = to[e];
        if (dis[u] + 1 == dis[v] && f[e]) {
            int ret = argument(v, min(rem - cnt,
               f[e]));
            f[e] -= ret;
            f[e ^ 1] += ret;
            cnt += ret;
        }
    }
    if (!cnt) dis[u] = -1;
    return cnt;
}



int dinic() {
    int cnt = 0, tmp;
    while (check())
    while (tmp = argument(S, INF))
    cnt += tmp;
    return cnt;
}
\end{verbatim}
\endgroup

\subsection{Luồng chi phí nhỏ nhất bằng SPFA}

\begingroup
\footnotesize
\begin{verbatim}
#include <bits/stdc++.h>
using namespace std;

const int N = 1100, INF = 0x3f3f3f3f;
const int M = N * N;

int pre[N], d[N], p[N], ans, cnt, hd[N], q[M], l, r;

struct edge {
    int u, v, w, c, nxt;
} e[M];

void init() {
    memset(hd, -1, sizeof(hd));
    ans = cnt = 0;
}

void add(int u, int v, int w, int c) {
    e[cnt].u = u, e[cnt].v = v, e[cnt].w = w,
       e[cnt].c = c, e[cnt].nxt = hd[u], hd[u] =
       cnt++;
    e[cnt].u = v, e[cnt].v = u, e[cnt].w = -w,
       e[cnt].c = 0, e[cnt].nxt = hd[v], hd[v] =
       cnt++;
}

void update(int s, int t) {
    int i, f = INF;
    for (i = t; i != s; i = e[pre[i]].u)
    f = min(f, e[pre[i]].c);


    for (i = t; i != s; i = e[pre[i]].u) {
        e[pre[i]].c -= f;
        e[pre[i] ^ 1].c += f;
        ans += f * e[pre[i]].w;
    }
}

int spfa(int s, int t) {
    int i, u, v, w;
    memset(p, 0, sizeof(p));
    memset(pre, -1, sizeof(pre));
    memset(d, 0x3f, sizeof(d));
    l = r = 0;
    q[++r] = s, p[s] = 1, d[s] = 0;
    while (l < r) {
        p[u = q[++l]] = 0;
        for (i = hd[u]; i != -1; i = e[i].nxt) {
            v = e[i].v, w = e[i].w;
            if (e[i].c && d[v] > d[u] + w) {
                d[v] = d[u] + w;
                pre[v] = i;
                if (!p[v]) {
                    p[v] = 1;
                    q[++r] = v;
                }
            }
        }
    }
    if (pre[t] == -1)
    return 0;
    return 1;
}
void solve(int s, int t) {
    ans = 0;
    while (spfa(s, t)) update(s, t);
}
\end{verbatim}
\endgroup

\subsection{Thuật toán Hungary}

\begingroup
\footnotesize
\begin{verbatim}
bool hungary(int u) {


    for (it = _v[u].begin(); it != _v[u].end(); it++)
       {
         int v = *it;
         if (!vis[v]) {
             vis[v] = true;
             if (!fa[v] || hungary(fa[v])) {
                 fa[v] = u;
                 return true;
             }
         }
    }
    return false;
}

int main() {
    for (int i = 1; i <= n; i++) {
        memset(vis, 0, sizeof vis);
        hungary(i);
    }
    return 0;
}
\end{verbatim}
\endgroup

\subsection{2-SAT}

\subsubsection{Kiểm tra tính thỏa được}

\begingroup
\footnotesize
\begin{verbatim}
#include <cstdio>
#include <algorithm>
#include <cstring>
#include <vector>
#include <stack>
#define N 10000
#define M 10000 + 5
using namespace std;

int T, n, m, x;
vector<int> V[N];
int main() {
    scanf("%d", &T);
    while (T--) {
        init();
        scanf("%d%d", &n, &m);
        for (int i = 1; i <= m; i++) {
            int a, b;
            char op1, op2;
            scanf("%*c%c%d %c%d", &op1, &a, &op2, &b);
            a <<= 1; if (op1 == 'h') a |= 1;
            b <<= 1; if (op2 == 'h') b |= 1;
            V[a ^ 1].push_back(b);
            V[b ^ 1].push_back(a);
        }
        for (int i = 2; i <= 2 * n + 1; i++) if
           (!dfn[i]) tarjan(i);
        bool f = true;
        for (int i = 1; i <= n; i++) {
            if (scc[i << 1] == scc[i << 1 | 1]) {
                f = false;
                break;
            }
        }
        if (f) puts("GOOD");
        else puts("BAD");
    }



    return 0;
}
\end{verbatim}
\endgroup

\subsubsection{Xuất một phương án}

\begingroup
\footnotesize
\begin{verbatim}
while(!q.empty()) {
    int u = q.front();
    q.pop();
    if(choose[u] != -1) continue;
    choose[u] = 0;
    choose[u^1] = 1;
    for(int i = 0; i < V[u].size(); i++) {
        --in[G2[u][i]];
        if(!in[V[u][i]]) q.push(V[u][i]);
    }
}
\end{verbatim}
\endgroup

Khi xuất phương án, hãy xây đồ thị ngược của đồ thị sau khi co các thành phần liên thông mạnh, rồi sắp xếp tô pô. Với mỗi đỉnh theo thứ tự tô pô (đỉnh này thực chất là một thành phần đã co), nếu nó chưa được gán màu thì chọn màu 1 cho nó và gán đỉnh đối ứng của nó là -1.

\subsection{Sắp xếp tô pô}

\begingroup
\footnotesize
\begin{verbatim}
#include <algorithm>
#include <cstring>
#include <cstdio>
#include <iostream>
#include <bitset>
#include <stack>
#define N (2000 + 5)
#define M (N * N >> 1) + 5
using namespace std;

int hd[N], nxt[M], to[M], tot;
int hd2[N], nxt2[M], to2[M], tot2;
int in[N];
void add(int a, int b) { nxt[++tot] = hd[a], to[hd[a]
   = tot] = b; }
void add2(int a, int b) { nxt2[++tot2] = hd2[a],
   to2[hd2[a] = tot2] = b, in[b]++; }


bitset<N> dp[N];
int n, _n, scc[N], siz[N], ans;

void topological_sort() {
    int q[N], head = 0, tail = -1;
    memset(q, 0, sizeof q);
    for (int i = 1; i <= _n; i++) {
        if (in[i] == 0) {
            q[++tail] = i;
            ans += siz[i] * dp[i].count();
        }
    }
    while (head <= tail) {
        int u = q[head++];
        for (int e = hd2[u]; e; e = nxt2[e]) {
            int v = to2[e];
            dp[v] |= dp[u];
            if (--in[v] == 0) {
                q[++tail] = v;
                ans += siz[v] * dp[v].count();
            }
        }
    }

}

bool vis[N];
int low[N], dfn[N], tim, x;
stack<int> s;
void tarjan(int u) {
    low[u] = dfn[u] = ++tim;
    s.push(u);
    vis[u] = true;
    for (int e = hd[u]; e; e = nxt[e]) {
        int v = to[e];
        if (!dfn[v]) {
            tarjan(v);
            low[u] = min(low[u], low[v]);
        } else if (vis[v]) low[u] = min(low[u],
           dfn[v]);


    }
    if (low[u] == dfn[u]) {
        _n++; do {
            x = s.top();
            s.pop();
            vis[x] = false;
            scc[x] = _n;
            siz[_n]++;
            dp[_n][x] = 1;
        } while (x != u);
    }
}

void prerequites() {
    for (int i = 1; i <= n; i++) if (!dfn[i])
       tarjan(i);
    for (int u = 1; u <= n; u++) {
        for (int e = hd[u]; e; e = nxt[e]) {
            int v = to[e];
            if (scc[v] != scc[u]) {
                add2(scc[u], scc[v]);
            }
        }
    }
}

int main() {
    scanf("%d", &n);
    for (int i = 1; i <= n; i++) {
        char s[N];
        scanf("%s", s + 1);
        for (int j = 1; j <= n; j++) {
            if (s[j] - '0') add(j, i);
        }
    }
    prerequites();
    topological_sort();
    printf("%d\n", ans);
    return 0;
}
\end{verbatim}
\endgroup

\subsection{Đường kính của cây}

\begingroup
\footnotesize
\begin{verbatim}
#include <cstdio>
#include <cstring>
#include <algorithm>
#include <cctype>
#include <cmath>
#define N 500000 + 5
#define inf 0x7fffffff
using namespace std;
#define mem(x) memset(x, 0, sizeof (x))

int fa[N], vis[N], dis[N], maxdis, maxi;
int hd[N], nxt[N], to[N], w[N], tot;
int n, s, diameter, ans;
int dia[N], dtot;
bool f;

void add(int a, int b, int c) {
    nxt[++tot] = hd[a], to[hd[a] = tot] = b, w[tot] =
       c;
    nxt[++tot] = hd[b], to[hd[b] = tot] = a, w[tot] =
       c;
}
void push(int u) { f = 1; while (u) { dia[++dtot] =
   u; u = fa[u]; } }
void dfs(int u, int Fa, int final) {
    if (final == 1 && dis[u] == maxdis && !f) {
       push(u); return; } fa[u] = Fa;
    for (int e = hd[u]; e; e = nxt[e]) {
        int v = to[e]; if (v != Fa) {
            if (final == 2 && vis[v]) continue;
            dis[v] = dis[u] + w[e];
            if (dis[v] > maxdis) maxdis = dis[v],
               maxi = v;
            dfs(v, u, final);
        }
    }
}
int main() {
    read(n), read(s); ans = inf;


    for (int i = 1; i < n; i++) { int a, b, c;
       read(a), read(b), read(c); add(a, b, c); }
    int tmp; dfs(1, 0, 0), maxdis = 0, mem(dis),
       mem(fa), dfs(tmp = maxi, 0, 0), dfs(tmp, 0, 1);
    diameter = maxdis;
    return 0;
}
\end{verbatim}
\endgroup

\subsection{Chia khối trên cây}

\begingroup
\footnotesize
\begin{verbatim}
#include <cstdio>
#include <algorithm>
#include <cstring>
#include <vector>
#define N 1000 + 5
using namespace std;
vector<int> _v[N];
int stack[N], top = -1;
int n, b;
int color[N], capital[N], cnt;

void dfs(int u, int fa) {
    int limit = top;
    vector<int>::iterator i;
    for (i = _v[u].begin(); i != _v[u].end(); i++) {
        int v = *i;
        if (v != fa) {
            dfs(v, u);
            if (top - limit >= b) {
                capital[++cnt] = u;
                while (top != limit)
                   color[stack[top--]] = cnt;
            }
        }
    }
    stack[++top] = u;
}

int main() {
    scanf("%d%d", &n, &b);


    for (int i = 1; i < n; i++) {
        int a, b;
        scanf("%d%d", &a, &b);
        _v[a].push_back(b);
        _v[b].push_back(a);
    }
    dfs(1, 0);
    while (top >= 0) color[stack[top--]] = cnt;
    printf("%d\n", cnt);
    for (int i = 1; i <= n; i++) printf("%d ",
       color[i]);
    puts("");
    for (int i = 1; i <= cnt; i++) printf("%d ",
       capital[i]);
    puts("");
    return 0;
}
\end{verbatim}
\endgroup

\subsection{Phân trị trọng tâm tĩnh}

\begingroup
\footnotesize
\begin{verbatim}
#include<cstdio>
#include<cstring>
#include<iostream>
#include<algorithm>
#define M 20200
using namespace std;
struct node {
    int to, f, nxt;
    bool ban;
} e[M << 1];
int head[M], tot = 1;
int n, ans;
int Siz[M];
void add(int x, int y, int z) {
    e[++tot].to = y;
    e[tot].f = z;
    e[tot].nxt = head[x];
    head[x] = tot;
}
void get_center(int x, int from, int Size, int &cg) {


   int i; Siz[x] = 1;
   bool flag = 1;
   for (i = head[x]; i; i = e[i].nxt) {
       if (e[i].ban || e[i].to == from) continue;
       get_center(e[i].to, x, Size, cg);
       if (Siz[e[i].to] > Size >> 1)
       flag = 0;
       Siz[x] += Siz[e[i].to];
   }
   if (Size - Siz[x] > Size >> 1) flag = 0;
   if (flag) cg = x;
}
void dfs(int x, int from, int dis, int cnt[]) {
    int i;
    Siz[x] = 1;
    cnt[dis]++;
    for (i = head[x]; i; i = e[i].nxt) {
        if (e[i].ban || e[i].to == from)
        continue;
        dfs(e[i].to, x, (dis + e[i].f) % 3, cnt);
        Siz[x] += Siz[e[i].to];
    }
}
void clac(int x) {
    int i;
    int cnt[3] = {1, 0, 0};
    for (i = head[x]; i; i = e[i].nxt) {
        if (e[i].ban) continue;
        int _cnt[3] = {0};
        dfs(e[i].to, x, e[i].f, _cnt);
        ans += cnt[0] * _cnt[0];
        ans += cnt[1] * _cnt[2];
        ans += cnt[2] * _cnt[1];
        cnt[0] += _cnt[0];
        cnt[1] += _cnt[1];
        cnt[2] += _cnt[2];
    }
}
void divide_and_conquer(int root, int Size) {
    int i, cg;
    if (Size == 1) return;


    get_center(root, 0, Size, cg);
    clac(cg);
    for (i = head[cg]; i; i = e[i].nxt) {
        if (e[i].ban)
        continue;
        e[i].ban = e[i ^ 1].ban = 1;
        divide_and_conquer(e[i].to, Siz[e[i].to]);
    }
}

int main() {
    int i, x, y, z;
    cin >> n;
    for (i = 1; i < n; i++) {
        scanf("%d%d%d", &x, &y, &z);
        add(x, y, z % 3);
        add(y, x, z % 3);
    }
    divide_and_conquer((1 + n) >> 1, n);
    ans = ans * 2 + n;
    int gcd = __gcd(ans, n * n);
    cout << ans / gcd << '/' << n*n / gcd << endl;
}
\end{verbatim}
\endgroup

\section{Cấu trúc dữ liệu}

\subsection{Cây Fenwick hai chiều}

\begingroup
\footnotesize
\begin{verbatim}
int lowbit(int x) {
    return x & -x;
}

void add(int x, int y, int d) {
    int i, j;
    for (i = x; i < N; i += lowbit(i))
    for (j = y; j < N; j += lowbit(j)) mat[i][j] += d;
}

int sum(int x, int y) {
    int ret = 0;
    int i, j;
    for (i = x; i > 0; i -= lowbit(i))
    for (j = y; j > 0; j -= lowbit(j)) ret +=
       mat[i][j];
    return ret;
}
\end{verbatim}
\endgroup

\subsection{Cây đoạn hai chiều}

\begingroup
\footnotesize
\begin{verbatim}
#include <bits/stdc++.h>
using namespace std;

const int maxn = 1000 * 1000 * 1.5;
// 2D line tree
struct node {
    short x1, y1, x2, y2; // (a, b) - (c, d)
    float max;
    int * ch;
};
node tree [maxn]; // space: O (N * M * 2) // In fact,
   1.4 * N * M is sufficient
int tol;
void maketree (int x1, int y1, int x2, int y2) {//
   time: O (2 * N * M)


    tol ++;
    int k = tol;
    tree [k] .x1 = x1;
    tree [k] .y1 = y1;
    tree [k] .x2 = x2;
    tree [k] .y2 = y2;
    tree [k] .max = 0.0;
    if (x1 == x2 && y1 == y2) {// c == d
        tree [k] .ch = 0;
        return;

    }
    tree [k] .ch = new int [4];
    int midx = (x1 + x2) / 2;
    int midy = (y1 + y2) / 2;
    if (x1 <= midx && y1 <= midy) {// lower left
        tree [k] .ch [0] = tol + 1;
        maketree (x1, y1, midx, midy);

    } else
    tree [k] .ch [0] = 0;
    if (midx + 1 <= x2 && midy + 1 <= y2) {// top
       right
        tree [k] .ch [1] = tol + 1;
        maketree (midx + 1, midy + 1, x2, y2);

    } else
    tree [k] .ch [1] = 0;
    if (x1 <= midx && midy + 1 <= y2) {// top left
        tree [k] .ch [2] = tol + 1;
        maketree (x1, midy + 1, midx, y2);

    } else
    tree [k] .ch [2] = 0;
    if (midx + 1 <= x2 && y1 <= midy) {// lower right
        tree [k] .ch [3] = tol + 1;
        maketree (midx + 1, y1, x2, midy);

    } else
    tree [k] .ch [3] = 0;
}


// time: O (logN) // N * N rectangle
void insert (int x, int y, int L, int k) {// Here is
   a rectangle for the insertion [x, y] - [x, y] .. A
   unit square
    if (tree [k] .x1 == tree [k] .x2 && tree [k] .y1
       == tree [k] .y2) {// leaf node //
        if (L > tree [k] .max)
        tree [k] .max = L;
        return;

    }
    int midx = (tree [k] .x1 + tree [k] .x2) / 2;
    int midy = (tree [k] .y1 + tree [k] .y2) / 2;
    if (x <= midx) {
        if (y <= midy)
        insert (x, y, L, tree [k] .ch [0]);
        else
        insert (x, y, L, tree [k] .ch [2]);

    } else {// x> midx
        if (y <= midy)
        insert (x, y, L, tree [k] .ch [3]);
        else
        insert (x, y, L, tree [k] .ch [1]);

    }
    float m = tree [tree [k] .ch [0]]. max;
    for (int i = 1; i < 4; i ++)
    m = max (m, tree [tree [k] .ch [i]]. max);
    tree [k] .max = m;
}
inline bool corss (int x1, int y1, int x2, int y2,
   int k) {// Determine if the two rectangles
   intersect
    int x3 = tree [k] .x1;
    int y3 = tree [k] .y1;
    int x4 = tree [k] .x2;
    int y4 = tree [k] .y2;
    if (y2 < y3 || y4 < y1 || x4 < x1 || x2 < x3)
    return false;
    else


    return true;
}
// time: O (logN) // N * N rectangle
float Query (int x1, int y1, int x2, int y2, int k) {
    if (corss (x1, y1, x2, y2, k) == false || tree
       [k] .max == 0) // The rectangles do not
       intersect or the rectangle max == 0 returns 0
    return 0;
    if (x1 <= tree [k] .x1 && y1 <= tree [k] .y1 &&
       tree [k] .x2 <= x2 && tree [k] .y2 <= y2) //
       If you want to query the rectangular coverage
       The current rectangle .. Then returns the
       current rectangle of the max value
    return tree [k] .max;
    // int midx = (tree [k] .x1 + tree [k] .x2) / 2;
    // int midy = (tree [k] .y1 + tree [k] .y2) / 2;
    float m [4] = {0, 0, 0, 0};
    for (int i = 0; i < 4; i ++)
    m [i] = Query (x1, y1, x2, y2, tree [k] .ch [i]);
    float mm = m [0];
    for (int i = 1; i < 4; i ++)
    mm = max (mm, m [i]);
    return mm;
}
\end{verbatim}
\endgroup

\subsection{Heap}

\begingroup
\footnotesize
\begin{verbatim}
#include <bits/stdc++.h>
/**
* This is a heap practice implement coded by
   Kvar_ispw17.
* email: enkerewpo@gmail.com
* blog: enkerewpo.github.io
*/
namespace __gnu_cxx_heap {

   int inline __left(int _x) {
       return _x * 2;
   }



int inline __right(int _x) {
    return _x * 2 + 1;
}

int inline __parent(int _x) {
    return floor(_x / 2);
}

template <typename _Tp> void inline
/**
* @brief This function will obtain the heap
   feature along in subheap rooted at @idx.
*        using the datas in array @A[] with size
   @heap_size.
* @param A         Input Array
* @param idx       Rooted being processed root
* @param heap_size the size of the full structure
*/
max_heapify(_Tp A[], int idx, size_t heap_size) {
    int l = __left(idx);
    int r = __right(idx);
    int max_e;
    if (l <= (int)heap_size && A[l] > A[idx])
       max_e = l;
    else max_e = idx;
    if (r <= (int)heap_size && A[r] > A[max_e])
       max_e = r;
    if (max_e != idx) {
        std::swap(A[max_e], A[idx]);
        max_heapify(A, max_e, heap_size);
    }
}

template <typename _Tp> void inline
/**
* @brief This is going to build a heap from
   pointer @begin to @end.
* @param begin pointer
* @param end   pointer
*/
build_max_heap(_Tp* begin, _Tp* end) {


    size_t size = end - begin - 1;
    for (int iter = size / 2; iter >= 0; iter--) {
        max_heapify(begin, iter, size);
    }
}

template <typename _Tp> inline
/**
* [__swap description]
* @param lhs [description]
* @param rhs [description]
*/
void __swap(_Tp &lhs, _Tp &rhs) {
    _Tp tmp = rhs;
    rhs = lhs;
    lhs = tmp;
}

template <typename _Tp> inline
/**
* sort on heap function
* @param begin pointer
* @param end   pointer
*/
void sort_heap(_Tp* begin, _Tp* end) {
    build_max_heap(begin, end);
    size_t size = end - begin - 1;
    for (int iter = size; iter >= 1; iter--) {
        __swap(begin[0], begin[iter]);
        max_heapify(begin, 0, --size);
    }
}

template <typename _Tp> inline
_Tp maximum_heap(_Tp* begin) {
    return begin[0];
}

template <typename _Tp> inline
/**
* pop heap top element to positoin @end


* @param begin pointer
* @param end   pointer
*/
void pop(_Tp* begin, _Tp* end) {
    size_t size = end - begin - 1;
    if (end - begin < 1) {
        fprintf(stderr, "heap size is zero when 
           it was instruct to pop()\n");
    }
    __swap(begin[0], begin[size]);
    size--;
    max_heapify(begin, 0, size);
}

template <typename _Tp> inline
/**
* increase key in heap
* @param A    array store _Tp vars
* @param idx index
* @param key keyword
*/
void heap_increase(_Tp A[], int idx, _Tp key) {
    A[idx] = key;
    while (idx > 1 and A[__parent(idx)] < A[idx])
       {
         __swap(A[idx], A[__parent(idx)]);
         idx = __parent(idx);
    }
}

template <typename _Tp> inline
/**
* push to heap function
* @param begin pointer
* @param end   pointer
* @param _x    input _Tp type varibles
*/
void push(_Tp* begin, _Tp* end, _Tp _x) {
    size_t size = end - begin;
    heap_increase(begin, size, _x);
}


}

int main() {
    int a[20] = {0, 4, 1, 3, 2, 16, 9, 10, 14, 8, 7};
    puts("original sequences:");
    for (int i = 1; i <= 10; i++) printf("%d ", a[i]);
    __gnu_cxx_heap::build_max_heap(a + 1, a + 10 + 1);
    printf("\n\n");

    puts("after build_max_heap:");
    for (int i = 1; i <= 10; i++) printf("%d ", a[i]);
    printf("\n\n");

    puts("after sort_heap:");
    __gnu_cxx_heap::sort_heap(a + 1, a + 10 + 1);
    for (int i = 1; i <= 10; i++) printf("%d ", a[i]);
    printf("\n\n");

    __gnu_cxx_heap::build_max_heap(a + 1, a + 10 + 1);
    puts("got maximum element:");
    printf("%d\n", __gnu_cxx_heap::maximum_heap(a +
       1));
    printf("\n");

    puts("after pop element:");
    __gnu_cxx_heap::pop(a + 1, a + 10 + 1);
    for (int i = 1; i <= 10; i++) printf("%d ", a[i]);
    printf("\n\n");

    __gnu_cxx_heap::build_max_heap(a + 1, a + 10 + 1);
    puts("after push 15:");
    __gnu_cxx_heap::push(a + 1, a + 11, 15);
    for (int i = 1; i <= 11; i++) printf("%d ", a[i]);
    printf("\n\n");

    return 0;
}
\end{verbatim}
\endgroup

\subsection{Heap trái}

\begingroup
\footnotesize
\begin{verbatim}
#include <bits/stdc++.h>
using namespace std;
/**
* This is a leftist heap practice implement coded by
   Kvar_ispw17.
* email: enkerewpo@gmail.com
* blog: enkerewpo.github.io
*/
class node {
    public:
    node* ls;
    node* rs;
    int x, npl;
};

typedef node heap;
typedef node* node_addr;

heap* mesh(heap* &lhs, heap* &rhs) {
    if (lhs == NULL) return rhs;
    if (rhs->x > lhs->x) swap(lhs, rhs);
    lhs->rs = mesh(lhs->rs, rhs);
    int l, r;
    if (lhs->ls == NULL) l = 0;
    else l = lhs->ls->npl;
    r = lhs->rs->npl;
    if (l < r) swap(lhs->ls, lhs->rs);
    else lhs->npl = lhs->rs->npl + 1;
    return lhs;
}

heap* meld(heap* &lhs, heap* &rhs) {
    if (lhs == NULL)
    return rhs;
    if (rhs == NULL)
    return lhs;
    return mesh(lhs, rhs);
}



void insert(int _x, heap* &h) {
    node* x = new node;
    x->x = _x;
    x->ls = NULL;
    x->rs = NULL;
    x->npl = 0;
    h = meld(x, h);
}

int pop(heap* &h) {
    int ret = h->x;
    h = meld(h->ls, h->rs);
    return ret;
}

int main() {
    heap* h = NULL;
    while (1) {
        int t;
        cin >> t;
        insert(t, h);
    }
    return 0;
}
\end{verbatim}
\endgroup

\subsection{Heap gộp được bằng PBDS}

\begingroup
\footnotesize
\begin{verbatim}
#include <ext/pb_ds/priority_queue.hpp>

using namespace __gnu_pbds;

typedef priority_queue<int, std::greater<int>,
   pairing_heap_tag, std::allocator<char> > heap;
heap p;
heap::iterator it;

int main() {
    p.push(5);
    p.push(2);
    p.push(10);


    p.push(1);
    for (int x : p) printf("%d ", x); // iterating by
       position
    printf("\n");
    for (it = p.begin(); it != p.end(); it++) {
        if (*it == 10) {
            p.erase(it);
        }
        if (*it == 2) {
            p.modify(it, 2333);
        }
    }
    for (int x : p) printf("%d ", x);
    printf("\n");
    heap tmp;
    tmp.push(142857);
    p.join(tmp);
    for (int x : p) printf("%d ", x);
    printf("\n");
    for (int x : tmp) printf("%d ", x);
    printf("\n");
    return 0;
}
\end{verbatim}
\endgroup

\subsection{Splay tree}

\begingroup
\footnotesize
\begin{verbatim}
#include <bits/stdc++.h>
#define MAXN 1000 + 5

int cnt, rt;
int Add[MAXN];

struct Tree {
    int key, num, size, fa, son[2];
} T[MAXN];

inline void PushUp (int x) {
    T[x].size = T[T[x].son[0]].size +
       T[T[x].son[1]].size + T[x].num;
}


inline void PushDown (int x) {
    if (Add[x]) {
        if (T[x].son[0]) {
            T[T[x].son[0]].key += Add[x];
            Add[T[x].son[0]] += Add[x];

        }
        if (T[x].son[1]) {
            T[T[x].son[1]].key += Add[x];
            Add[T[x].son[1]] += Add[x];

        }
        Add[x] = 0;

    }
}

inline int Newnode (int key, int fa) {// Create a new
   node and return
    ++ cnt;
    T[cnt].key = key;
    T[cnt].num = T[cnt].size = 1;
    T[cnt].fa = fa;
    T[cnt].son[0] = T[cnt].son[1] = 0;
    return cnt;
}

inline void Rotate (int x, int p) {// 0 L-1
   dextrorotation
    int y = T[x].fa;
    PushDown (y);
    PushDown (x);
    T[y].son[!p] = T[x].son[p];
    T[T[x].son[p]].fa = y;
    T[x].fa = T[y].fa;
    if (T[x].fa)
    T[T[x].fa].son[T[T[x].fa].son[1] == y] = x;
    T[x].son[p] = y;
    T[y].fa = x;
    PushUp (y);


    PushUp (x);
}

void Splay (int x, int To) {// Move the x node to the
   child of To
    while (T[x].fa != To) {
        if (T[T[x].fa].fa == To)
        Rotate (x, T[T[x].fa].son[0] == x);
        else {
            int y = T[x].fa, z = T[y].fa;
            int p = (T[z].son[0] == y);
            if (T[y].son[p] == x) Rotate (x, ! p),
               Rotate (x, p);
            else Rotate (y, p), Rotate (x, p);// a
               word spin

        }

    }
    if (To == 0) rt = x;
}

int GetPth (int p, int To) {// Returns the pth
   smallest node and moves to the child of To
    if (! rt || p > T[rt].size) return 0;
    int x = rt;
    while (x) {
        PushDown (x);
        if (p >= T[T[x].son[0]].size + 1 && p <=
           T[T[x].son[0]].size + T[x].num)
        break;
        if (p > T[T[x].son[0]].size + T[x].num) {
            p -= T[T[x].son[0]].size + T[x].num;
            x = T[x].son[1];

        } else
        x = T[x].son[0];

    }
    Splay (x, 0);
    return x;


}

int Find (int key) {
    // Returns the key value of the node if not
       returned 0 If it is transferred to the root
    if (! rt) return 0;
    int x = rt;
    while (x) {
        PushDown (x);
        if (T[x].key == key) break;
        x = T[x].son[key > T[x].key];

    }
    if (x) Splay (x, 0);
    return x;
}

int Prev () {
    // return to the root node of the non - focus
    if (! rt || ! T[rt].son[0]) return 0;
    int x = T[rt].son[0];
    while (T[x].son[1]) {
        PushDown (x);
        x = T[x].son[1];

    }
    Splay (x, 0);
    return x;
}

int Succ () {// Returns the successor to the root node
    if (! rt || ! T[rt].son[1]) return 0;
    int x = T[rt].son[1];
    while (T[x].son[0]) {
        PushDown (x);
        x = T[x].son[0];

    }
    Splay (x, 0);
    return x;
}


void Insert (int key) {// Insert the key value
    if (! rt)
    rt = Newnode (key, 0);
    else {
        int x = rt, y = 0;
        while (x) {
            PushDown (x);
            y = x;
            if (T[x].key == key) {
                T[x].num ++;
                T[x].size ++;
                break;

            }
            T[x].size ++;
            x = T[x].son[key > T[x].key];

        }
        if (! x)
        x = T[y].son[key > T[y].key] = Newnode (key,
           y);
        Splay (x, 0);

    }
}

void Delete (int key) {// Delete 1 node whose value
   is key
    int x = Find (key);
    if (! x) return;
    if (T[x].num > 1) {
        T[x].num--;
        PushUp (x);
        return;

    }
    int y = T[x].son[0];
    while (T[y].son[1])
    y = T[y].son[1];
    int z = T[x].son[1];


    while (T[z].son[0])
    z = T[z].son[0];
    if (! y && ! z) {
        rt = 0;
        return;

    }
    if (! y) {
        Splay (z, 0);
        T[z].son[0] = 0;
        PushUp (z);
        return;

    }
    if (! z) {
        Splay (y, 0);
        T[y].son[1] = 0;
        PushUp (y);
        return;

    }
    Splay (y, 0);
    Splay (z, y);
    T[z].son[0] = 0;
    PushUp (z);
    PushUp (y);
}

int GetRank (int key) {// Get the number of nodes
   with value <= key
    if (! Find (key)) {
        Insert (key);
        int tmp = T[T[rt].son[0]].size;
        Delete (key);
        return tmp;

    } else
    return T[T[rt].son[0]].size + T[rt].num;
}




void Delete (int l, int r) {// Delete all nodes with
   values   in[l, r] l!= r
    if (! Find (l)) Insert (l);
    int p = Prev ();
    if (! Find (r)) Insert (r);
    int q = Succ ();
    if (! p && ! q) {
        rt = 0;
        return;

    }
    if (! p) {
        T[rt].son[0] = 0;
        PushUp (rt);
        return;

    }
    if (! q) {
        Splay (p, 0);
        T[rt].son[1] = 0;
        PushUp (rt);
        return;

    }
    Splay (p, q);
    T[p].son[1] = 0;
    PushUp (p);
    PushUp (q);
}
\end{verbatim}
\endgroup

\subsection{Splay tree bằng PBDS}

\begingroup
\footnotesize
\begin{verbatim}
#include <ext/pb_ds/assoc_container.hpp>
#include <ext/pb_ds/tree_policy.hpp>
#include <bits/stdc++.h>
using namespace __gnu_pbds;
using namespace std;
typedef tree<int, null_type, less<int>,
   splay_tree_tag, tree_order_statistics_node_update>
   splay;


splay s;
splay :: iterator it;
int main() {

    for (int i = 20; i >= 15; i--) s.insert(i); //15
       16 17 18 19 20
    for (it = s.begin(); it != s.end(); it++)
       printf("%d ", *it);
    printf("\n");
    it = s.find_by_order(3);
    printf("NO.4 = %d\n", *it);      // return NO.k+1
       if you input find_by_order(k)
    int rnk = s.order_of_key(16);
    printf("RANK OF 16 is %d\n", rnk + 1); // return
       how many items are "less" than your input
    splay t; t.clear();
    t.join(s); // copy 's' to 't'
    printf("    t = ");
    for (it = t.begin(); it != t.end(); it++)
       printf("%d ", *it);
    printf("\n");
    splay u; u.clear();
    t.split(17, u); // clear 'u' and move items
       "greater" then 17 to 'u'
    printf("    u = ");
    for (it = u.begin(); it != u.end(); it++)
       printf("%d ", *it);
    printf("\n");
    return 0;
}
\end{verbatim}
\endgroup

\subsection{bitset}

\begingroup
\footnotesize
\begin{verbatim}
#include <cstdio>
#include <iostream>
#include <bitset>
using namespace std;
#define MAXN 20

int main() {


    bitset<5> a("10010");
    bitset<5> b("01110");
    cout << (a ^ b) << endl; // 11100
    cout << (a & b) << endl; // 00010
    a.flip();
    cout << a << endl; // 01101
    return 0;
}
\end{verbatim}
\endgroup

\subsection{Cây đoạn bền vững}

\begingroup
\footnotesize
\begin{verbatim}
#include<cstdio>
#include<cstring>
#include<iostream>
#include<algorithm>
#include<vector>

using namespace std;
const int MAXN = 1e5 + 5;
int n, m, cnt, root[MAXN], a[MAXN], x, y, k;

struct node {
    int l, r, sum;
} T[MAXN * 30];
vector<int> v;

int getid(int x) {
    return lower_bound(v.begin(), v.end(), x) -
       v.begin() + 1;
}

void update(int l, int r, int &x, int y, int pos) {
    T[++cnt] = T[y], T[cnt].sum++, x = cnt;
    if (l == r) return;
    int mid = (l + r) >> 1;
    if (mid >= pos) update(l, mid, T[x].l, T[y].l,
       pos);
    else update(mid + 1, r, T[x].r, T[y].r, pos);
}



int query(int l, int r, int x, int y, int k) {
    if (l == r) return l;
    int mid = (l + r) >> 1;
    int sum = T[T[y].l].sum - T[T[x].l].sum;
    if (sum >= k) return query(l, mid, T[x].l,
       T[y].l, k);
    else return query(mid + 1, r, T[x].r, T[y].r, k -
       sum);
}

int main() {
    while (scanf("%d%d", &n, &m) != EOF) {
    v.clear();
    memset(root, 0, sizeof(root));
    for (int i = 1; i <= n; i++) {
        scanf("%d", &a[i]);
        v.push_back(a[i]);
    }
    sort(v.begin(), v.end());
    v.erase(unique(v.begin(), v.end()), v.end());
    for (int i = 1; i <= n; i++) {
        update(1, n, root[i], root[i - 1],
           getid(a[i]));
    }
    while (m--) {
            scanf("%d%d%d", &x, &y, &k);
            printf("%d\n", v[query(1, n, root[x - 1],
               root[y], k) - 1]);
        }
    }
    return 0;
}
\end{verbatim}
\endgroup

\subsection{Skip list}

\begingroup
\footnotesize
\begin{verbatim}
#include <cstdio>

#define MAX_LEVEL 10 // The maximum number of layers

//node
typedef struct nodeStructure {
    int key;
    int value;
    struct nodeStructure *forward[1];
} nodeStructure;

// jump table
typedef struct skiplist {
    int level;
    nodeStructure *header;
} skiplist;

// Create the node
nodeStructure *createNode(int level, int key, int
   value)
{
    nodeStructure *ns =(nodeStructure *)
       malloc(sizeof(nodeStructure) + level
       *sizeof(nodeStructure *));
    ns->key = key;
    ns->value = value;
    return ns;
}

// Initialize the jump table
skiplist *createSkiplist()
{
    skiplist *sl =(skiplist *)
       malloc(sizeof(skiplist));
    sl->level = 0;
    sl->header = createNode(MAX_LEVEL-1,0,0);
    for(int i = 0; i <MAX_LEVEL; i++) {
        sl->header->forward[i] = NULL;
    }


    return sl;
}

// Randomly generated layers
int randomLevel()
{
    int k = 1;
    while(rand()% 2)
    k++;
    k =(k <MAX_LEVEL)? k: MAX_LEVEL;
    return k;
}

// insert node
bool insert(skiplist *sl, int key, int value)
{
    nodeStructure *update[MAX_LEVEL];
    nodeStructure *p, *q = NULL;
    p = sl->header;
    int k = sl->level;
    // Find the location you want to insert from the
       top
    // fill update
    for(int i = k-1; i >= 0; i--) {
        while((q = p->forward[i]) &&(q->key <key)) {
            p = q;
        }
        update[i] = p;
    }
    // can not insert the same key
    if(q && q->key == key) {
        return false;
    }

    // Generate a random number of layers K
    // Create a new node to insert q
    // Insert layer by layer
    k = randomLevel();
    // update the level of the jump table
    if(k>(sl->level)) {
        for(int i = sl->level; i <k; i++) {


            update[i] = sl->header;
        }
        sl->level = k;
    }

    q = createNode(k, key, value);
    // Update node pointer layer by layer, same as
       normal list insertion
    for(int i = 0; i <k; i++) {
        q->forward[i] = update[i]->forward[i];
        update[i]->forward[i] = q;
    }
    return true;
}

// Search the value of the specified key
int search(skiplist *sl, int key)
{
    nodeStructure *p, *q = NULL;
    p = sl->header;
    // Search from the highest level
    int k = sl->level;
    for(int i = k-1; i >= 0; i--) {
        while((q = p->forward[i]) &&(q->key <= key)) {
            if(q->key == key) {
                return q->value;
            }
            p = q;
        }
    }
    return NULL;
}

// delete the specified key
bool deleteSL(skiplist *sl, int key)
{
    nodeStructure *update[MAX_LEVEL];
    nodeStructure *p, *q = NULL;
    p = sl->header;
    // Search from the highest level
    int k = sl->level;


    for(int i = k-1; i >= 0; i--) {
        while((q = p->forward[i]) &&(q->key <key)) {
            p = q;
        }
        update[i] = p;
    }
    if(q && q->key == key) {
        / / Delete layer by layer, and delete the
           same general list
        for(int i = 0; i <sl->level; i++) {
            if(update[i]->forward[i] == q) {
                 update[i]->forward[i] = q->forward[i];
            }
        }
        free(q);
        // If you delete the node is the largest
           layer, you need to re-maintain the jump
           table
        for(int i = sl->level - 1; i >= 0; i--) {
            if(sl->header->forward[i] == NULL) {
                 sl->level--;
            }
        }
        return true;
    } else
    return false;
}

void printSL(skiplist *sl)
{
    // print from the highest level
    nodeStructure *p, *q = NULL;

    // Search from the highest level
    int k = sl->level;
    for(int i = k-1; i >= 0; i--) {
        p = sl->header;
        while(q = p->forward[i]) {
            printf("% d->", p->value);
            p = q;
        }


       printf("\n");
   }
   printf("\n");
}
int main()
{
    skiplist *sl = createSkiplist();
    for(int i = 1; i <= 19; i++) {
        insert(sl, i, i *2);
    }
    printSL(sl);
    //search for
    int i = search(sl, 4);
    printf("i =% d \n", i);
    //delete
    bool b = deleteSL(sl, 4);
    if(b) printf("delete successfully \n");
    printSL(sl);
    return 0;
}
\end{verbatim}
\endgroup

\section{Quy hoạch động}

\subsection{Tối ưu hóa bằng độ dốc}

\begingroup
\footnotesize
\begin{verbatim}
// BZOJ 1991
#include <cstdio>
#include <cstring>
#include <cctype>
#include <algorithm>
#include <cmath>
#define slope_optimize
using namespace std;
#define N 1000000 + 5
typedef long long ll;
#define int ll
int n;
ll a, b, c;
ll A[N], S[N], dp[N];
struct point {
    point() {}
    point(int _, int __) : x(_), y(__) {}
    point(int _, int __, int ___) : x(_), y(__),
       id(___) {}
    int x, y, id;
};
template <typename _Tp>
void read(_Tp &a, char c = 0, int op = 1) {
    for (c = getchar(); !isdigit(c); c = getchar()) c
       == '-' ? op = -1 : op = 1;
    for (a = 0; isdigit(c); a = a * 10 + c - '0', c =
       getchar());
    a *= op;
}
double slp(point a, point b) { return double(b.y -
   a.y) / double(b.x - a.x); }
ll w(int x, int y) { ll s = S[y] - S[x - 1]; return a
   * s * s + b * s + c; }
point q[N];
int l, r;
signed main() {



    read(n), read(a), read(b), read(c);
    for (int i = 1; i <= n; i++) read(A[i]), S[i] =
       S[i - 1] + A[i];
    slope_optimize {
        for (int i = 1; i <= n; i++) {
            double k = S[i];
            while (l < r && slp(q[l], q[l + 1]) < k)
               l++;
            point bst = q[l];
            dp[i] = dp[bst.id] + w(bst.id + 1, i);
            point now = point(2 * a * S[i], dp[i] + a
               * S[i] * S[i] - b * S[i], i);
            while (l < r && slp(q[r], now) < slp(q[r
               - 1], q[r])) r--;
            q[++r] = now;
        }
    }
    printf("%lld\n", dp[n]);
    return 0;

}
\end{verbatim}
\endgroup

\subsection{Bài toán ba lô có phụ thuộc}

\begingroup
\footnotesize
\begin{verbatim}
#include <cstdio>
#include <algorithm>
#include <stack>
#include <cstring>

#define N 100 + 5
#define M 500 + 5
using namespace std;

int dp[N][M];

int hd[N], nxt[N], to[N], tot;
int hd2[N], nxt2[N], to2[N], in2[N], tot2;
void add (int a, int b) { nxt [++tot ] = hd [a], to
   [hd [a] = tot ] = b; }



void add2(int a, int b) { nxt2[++tot2] = hd2[a],
   to2[hd2[a] = tot2] = b, in2[b]++; }

int n, m, _n, color[N], low[N], dfn[N], tim, sv[N],
   sw[N];
int v[N], w[N];
bool vis[N];
stack<int> s;

void tarjan(int u) {
    low[u] = dfn[u] = ++tim;
    s.push(u);
    vis[u] = true;
    for (int e = hd[u]; e; e = nxt[e]) {
        int v = to[e];
        if (!dfn[v]) {
            tarjan(v);
            low[u] = min(low[u], low[v]);
        } else if (vis[v])low[u] = min(low[u],
           dfn[v]);
    }
    if (low[u] == dfn[u]) {
        _n++; int x; do {
            x = s.top();
            s.pop();
            vis[x] = false;
            color[x] = _n;
            sv[_n] += v[x];
            sw[_n] += w[x];
        } while (x != u);
    }
}

void dfs(int i) {
    for (int e = hd2[i]; e; e = nxt2[e]) {
        int v = to2[e];
        dfs(v);
        for (int j = m - sw[i]; j >= 0; j--) {
            for (int k = 0; k <= j; k++) {
                dp[i][j] = max(dp[i][j], dp[i][k] +
                   dp[v][j - k]);


            }
        }
    }
    for (int j = m; j >= 0; j--) {
        if (j >= sw[i]) dp[i][j] = dp[i][j - sw[i]] +
           sv[i];
        else dp[i][j] = 0;
    }
}

int main() {
    scanf("%d%d", &n, &m);
    for (int i = 1; i <= n; i++) scanf("%d", &w[i]);
    for (int i = 1; i <= n; i++) scanf("%d", &v[i]);
    for (int i = 1; i <= n; i++) {
        int a;
        scanf("%d", &a);
        if (a) add(a, i);
    }
    for (int i = 1; i <= n; i++) if (!dfn[i])
       tarjan(i);
    for (int i = 1; i <= n; i++) {
        for (int e = hd[i]; e; e = nxt[e]) {
            int v = to[e];
            if (color[v] != color[i]) add2(color[i],
               color[v]);
        }
    }
    for (int i = 1; i <= _n; i++) if (in2[i] == 0)
       add2(_n + 1, i);
    dfs(_n + 1);
    printf("%d\n", dp[_n + 1][m]);
    return 0;
}
\end{verbatim}
\endgroup

\subsection{Dãy con tăng chung dài nhất}

\begingroup
\footnotesize
\begin{verbatim}
#include <bits/stdc++.h>

using namespace std;


const int N = 5005;

int a[N], b[N], f[N][N], g[N][N];

void work(int x, int y, int z) {
    if (!z) return;
    while (a[x] != b[y]) --x;
    work(x, g[x][y], z - 1);
    printf("%d ", b[y]);
}
int main() {
    int n, m, i, tma, pre, j, ans = 0, t;
    scanf("%d", &n);
    for (i = 1; i <= n; ++i) scanf("%d", &a[i]);
    scanf("%d", &m);
    for (i = 1; i <= m; ++i) scanf("%d", &b[i]);
    for (i = 1; i <= n; ++i)
    for (tma = 0, j = 1; j <= m; ++j) {
        if (b[j] < a[i] && f[i - 1][j] > tma) tma =
           f[i - 1][pre = j];
        if (a[i] == b[j]) {f[i][j] = tma + 1;
           g[i][j]=pre;} else f[i][j] = f[i - 1][j];
        if (f[i][j] > ans) ans = f[i][t = j];
    }
    printf("%d\n", ans);
    if (ans) {
        work(n, t, ans);
        printf("\n");
    }
    return 0;
}
\end{verbatim}
\endgroup

\section{Khác}

\subsection{Thuật toán Mo có cập nhật}

\begingroup
\footnotesize
\begin{verbatim}
#include <cstdio>
#include <algorithm>
#include <cstring>
#include <cmath>
#define N 1000000 + 5
#define M 1000000 + 5
using namespace std;
#define int long long
int n, m, block;
int a[N], tmp[N], ans[M];

int qcnt;
struct query {
    query() {}
    query(int _l, int _r, int _d, int _i) : l(_l),
       r(_r), dfn(_d), id(_i) {}
    int l, r, dfn, id;
    bool operator < (const query &rhs) const {
        return (l / block) == (rhs.l / block) ? (r ==
           rhs.r ? dfn < rhs.dfn : r < rhs.r) : l /
           block < rhs.l / block;
    }

} q[M];


int ccnt;
struct change {
    change() {}
    change(int _p, int _v, int _pr) : pos(_p),
       val(_v), pre(_pr) {}
    int pos, val, pre;

} c[M];

int cnt[N], tot;



int L, R, DFN;

void add(int p) { tot += (++cnt[p]) == 1; }
void del(int p) { tot -= (--cnt[p]) == 0; }
void rev(int p, int v) {
    if(L <= p && p <= R) add(v), del(a[p]);
    a[p] = v;
}

void mo() {
    for(int i = 1; i <= qcnt; i++) {
        while(DFN < q[i].dfn) DFN++, rev(c[DFN].pos,
           c[DFN].val);
        while(DFN > q[i].dfn) rev(c[DFN].pos,
           c[DFN].pre), DFN--;
        while(R < q[i].r) R++, add(a[R]);
        while(R > q[i].r) del(a[R]), R--;
        while(L < q[i].l) del(a[L]), L++;
        while(L > q[i].l) L--, add(a[L]);
        ans[q[i].id] = tot;
    }
}

signed main() {
    scanf("%lld%lld", &n, &m);
    block = sqrt(n);
    for(int i = 1; i <= n; i++) {
        scanf("%lld", &a[i]);
        tmp[i] = a[i];
    }
    int a, b;
    for(int i = 1; i <= m; i++) {
        char op[3];
        scanf("%s%lld%lld", op, &a, &b);
        if(op[0] == 'Q') qcnt++, q[qcnt] = query(a,
           b, ccnt, qcnt);
        else c[++ccnt] = change(a, b, tmp[a]), tmp[a]
           = b;
    }
    sort(q + 1,q + qcnt + 1);
    mo();


    for(int i = 1; i <= qcnt; i++) printf("%lld\n",
       ans[i]);
    return 0;
}
\end{verbatim}
\endgroup

\end{document}
